\documentclass[11pt]{article}

\usepackage[a4paper,margin=1in]{geometry}
\usepackage{amsmath,amssymb,mathtools}
\usepackage{booktabs}
\usepackage{tabularx}
\usepackage{enumitem}
\usepackage[T1]{fontenc}
\usepackage{lmodern}

\title{An Abstract MILP Formulation for Discrete Pallet Packing}
\author{}
\date{}

\newcommand{\Item}{\mathcal{I}}
\newcommand{\Pallet}{\mathcal{P}}
\newcommand{\Place}{\mathcal{Q}}
\newcommand{\Cell}{\mathcal{C}}
\newcommand{\Support}{\mathcal{R}}
\newcommand{\Prio}{\pi}

\begin{document}

\maketitle

\begin{abstract}
This document gives an abstract mixed-integer linear programming (MILP)
formulation for the discrete pallet-packing model implemented in the project.
The formulation describes the mathematical structure of the model rather than
the instance-specific expansion generated by an optimizer.  Candidate box
placements are enumerated before optimization; the MILP then chooses exactly
one candidate placement for each item while enforcing pallet capacity,
non-overlap, support, compatibility, accessibility, and symmetry conditions.
Several linear objectives can be optimized sequentially to obtain a Pareto
front between compactness, accessibility, and vertical center-of-mass quality.
\end{abstract}

\section{Modeling approach}

The physical dimensions are discretized using a grid of size $g$ millimetres.
For each item, allowable orientations and feasible grid coordinates are
enumerated in a preprocessing step.  A \emph{placement} is therefore a fully
specified choice of item, pallet, orientation, and three-dimensional position.
The MILP does not search over continuous coordinates: it selects from this
finite placement set.

This position-indexed construction has two useful consequences.  First,
geometric coefficients such as overlap areas and occupied cells are constants
known before optimization.  Second, non-overlap and support can be expressed
with linear inequalities over binary placement variables.

\section{Sets and indices}

\begin{description}[leftmargin=2.4cm,style=nextline]
  \item[$i,j\in\Item$] Items (individual boxes) that must be loaded.
  \item[$p\in\Pallet$] Candidate pallet positions, indexed from $1$ to
    $P^{\max}$.  The model may use fewer than $P^{\max}$ pallets.
  \item[$q\in\Place$] All feasible candidate placements generated during
    preprocessing.  Each placement belongs to exactly one item and one
    candidate pallet.
  \item[$\Place_i\subseteq\Place$] Placements of item $i$.
  \item[$\Place_p\subseteq\Place$] Placements assigned to candidate pallet $p$.
  \item[$\Place_{ip}$] Placements of item $i$ on pallet $p$.
  \item[$\Place(c)\subseteq\Place$] Placements that occupy grid cell
    $c\in\Cell$.  The cells are the discretized locations inside all pallets.
  \item[$\Support(q)\subseteq\Place$] Candidate placements that can directly
    support placement $q$: they are on the same pallet, their top is exactly
    at the base of $q$, and their footprints overlap.
  \item[$\mathcal{G}$] Groups of physically identical items.  Each group is
    used for symmetry-breaking constraints.
  \item[$\mathcal{H}$] Hard incompatibility relations.  An element
    $(B,U)\in\mathcal{H}$ contains a set $B$ of possible lower placements and
    a set $U$ of possible upper placements that may not be selected together.
    These relations are generated from food--chemical and fragile-item rules.
  \item[$\mathcal{A}$] Ordered item pairs $(i,j)$ for which item $i$ has higher
    retrieval priority than item $j$.
  \item[$\mathcal{R}_{ij}$] Column relations $(B,U)$ that constitute a
    possible accessibility inversion for priority pair $(i,j)\in\mathcal{A}$.
\end{description}

\section{Input data and derived placement data}

The following data are read from the instance and configuration files:

\begin{description}[leftmargin=2.4cm,style=nextline]
  \item[$L,W,H$] Pallet length, width, and height in grid units.
  \item[$K$] Required number of pallets for the current optimization run.
  \item[$d_i^x,d_i^y,d_i^z$] Dimensions of item $i$ after conversion to grid
    units.  The dimensions are rounded up to preserve physical feasibility.
  \item[$w_i$] Weight of item $i$.
  \item[$\rho_i$] Density of item $i$ when the density-weighted vertical
    objective is selected.
  \item[$\Prio_i$] Retrieval priority of item $i$.
  \item[$f_i,c_i,\varphi_i$] Binary item attributes indicating food,
    chemical, and fragile status, respectively.
  \item[$z_q,d_q^z$] Base height and vertical size of placement $q$.
  \item[$t_q=z_q+d_q^z$] Top height of placement $q$.
  \item[$A_q=d_q^x d_q^y$] Base area of placement $q$.
  \item[$A_{qr}$] Footprint-overlap area between placements $q$ and $r$.
  \item[$\alpha$] Minimum support fraction, with $0\le\alpha\le1$.  The
    default configuration uses $\alpha=0.75$.
\end{description}

For every candidate placement $q$, preprocessing also stores its item index
$i(q)$, pallet index $p(q)$, orientation, grid coordinates, and occupied cells.
All of these values are constants in the MILP.  For symmetry breaking,
$r_{iq}$ denotes the fixed rank of placement $q$ within the ordered placement
list of item $i$.

\section{Decision variables}

\begin{alignat}{3}
  x_q &\in\{0,1\} && && \text{1 if placement $q$ is selected,}\\
  y_p &\in\{0,1\} && && \text{1 if candidate pallet $p$ is used,}\\
  h_p &\in[0,H] && && \text{maximum occupied height on pallet $p$,}\\
  H^{\max},H^{\min},s &\in[0,H] && && \text{maximum height, minimum height, and spread,}\\
  a_{ij} &\in\{0,1\} && && \text{1 if the priority pair $(i,j)$ violates access order.}
\end{alignat}

The placement variables $x_q$ are the main combinatorial variables.  The
continuous height variables summarize the selected placements, while $a_{ij}$
are soft-penalty variables used only when accessibility is configured as a
soft objective.

\section{Core feasibility constraints}

\subsection{Exactly one placement per item}

Every item must be placed exactly once:
\begin{equation}
  \sum_{q\in\Place_i}x_q = 1
  \qquad \forall i\in\Item.
  \label{eq:place-once}
\end{equation}

\subsection{Pallet assignment and pallet count}

A selected placement on pallet $p$ activates that pallet.  Conversely, an item
assigned to $p$ may only use $p$ when it is active:
\begin{align}
  \sum_{q\in\Place_p}x_q &\ge y_p && \forall p\in\Pallet, \\
  \sum_{q\in\Place_{ip}}x_q &\le y_p && \forall i\in\Item,\ p\in\Pallet, \\
  \sum_{p\in\Pallet}y_p &= K. &&
  \label{eq:pallet-assignment}
\end{align}

The candidate pallets are symmetry-broken by requiring earlier pallet indices
to be used before later ones:
\begin{equation}
  y_p\ge y_{p+1}
  \qquad p=1,\ldots,P^{\max}-1.
  \label{eq:pallet-order}
\end{equation}

\subsection{Stack height and height spread}

If placement $q$ is selected, the height of its pallet reaches its top:
\begin{equation}
  h_{p(q)} \ge t_q x_q
  \qquad \forall q\in\Place.
  \label{eq:height-placement}
\end{equation}

An unused pallet has zero height, and the global extrema summarize the pallet
heights:
\begin{align}
  h_p &\le H y_p && \forall p\in\Pallet,\\
  H^{\max} &\ge h_p && \forall p\in\Pallet,\\
  H^{\min} &\le h_p + H(1-y_p) && \forall p\in\Pallet,\\
  s &= H^{\max}-H^{\min}. &&
  \label{eq:height-spread}
\end{align}

\subsection{Geometric non-overlap}

For every discretized cell, at most one selected placement may occupy that
cell:
\begin{equation}
  \sum_{q\in\Place(c)}x_q\le 1
  \qquad \forall c\in\Cell.
  \label{eq:nonoverlap}
\end{equation}

This cell-clique representation is equivalent to forbidding volumetric overlap
among selected placements, while avoiding a separate binary variable for every
pair of placements.

\subsection{Support of boxes above the floor}

If support constraints are enabled, every selected placement above the floor
must receive at least an $\alpha$ fraction of its base area from placements
directly below it:
\begin{equation}
  \sum_{r\in\Support(q)} A_{qr}x_r
  \ge \alpha A_q x_q
  \qquad \forall q\in\Place:\ z_q>0.
  \label{eq:support}
\end{equation}

The coefficients $A_{qr}$ are known overlap areas, so the constraint is
linear.  Placements with $z_q=0$ require no supporter because they rest on the
pallet floor.

\section{Compatibility and handling constraints}

\subsection{Food--chemical separation}

When food and chemical overlap is prohibited, a food placement and a chemical
placement cannot occupy the same vertical column with one below the other.
For every such column and height threshold, let $B$ be the compatible lower
placements and $U$ the incompatible upper placements.  The model imposes:
\begin{equation}
  \sum_{q\in B}x_q + \sum_{q\in U}x_q \le 1.
  \label{eq:food-chemical}
\end{equation}

The same constraint is generated in both vertical directions, so a food item
cannot be above or below a chemical item when their footprints overlap.

\subsection{Fragile items cannot support other items}

For every fragile lower item and every possible item above it, the same column
conflict construction gives:
\begin{equation}
  \sum_{q\in B_{\mathrm{fragile}}}x_q
  +\sum_{q\in U}x_q \le 1.
  \label{eq:fragile}
\end{equation}

Thus a fragile selected placement cannot be a direct supporter of another
selected placement.

\subsection{Symmetry breaking for identical items}

If two consecutive items in an identical-item group have placement lists
ordered as $\Place_i=(q_1,\ldots,q_m)$ and
$\Place_j=(r_1,\ldots,r_m)$, the selected placement ranks are ordered:
\begin{equation}
  \sum_{k=1}^{m} kx_{q_k}
  \le
  \sum_{k=1}^{m} kx_{r_k}.
  \label{eq:symmetry}
\end{equation}

This does not change the set of physically distinguishable packings; it only
removes equivalent permutations of identical boxes.

\section{Accessibility objective and constraints}

Suppose item $i$ has higher retrieval priority than item $j$, so that $i$
should not be blocked by $j$.  For each relevant vertical-column relation, the
binary violation variable satisfies:
\begin{equation}
  a_{ij}\ge
  \sum_{q\in B_{ij}}x_q + \sum_{q\in U_{ij}}x_q - 1.
  \label{eq:access-violation}
\end{equation}

The variable can remain zero when no blocking relation is selected and must be
one if both sides of a forbidden priority relation are selected.  The total
number of accessibility inversions is:
\begin{equation}
  A=\sum_{(i,j)}a_{ij}.
  \label{eq:access-objective}
\end{equation}

If accessibility is a hard requirement, the corresponding right-hand side is
replaced by zero-violation inequalities.  In the default soft formulation,
$A$ is minimized as a secondary objective.

\section{Vertical moment objectives}

The project supports two alternative linear measures of vertical loading.
For a density-weighted objective:
\begin{equation}
  M^{\mathrm{density}}
  =\sum_{q\in\Place}
  \rho_{i(q)}
  \left(\frac{g(z_q+d_q^z/2)}{1000}\right)x_q.
  \label{eq:density-moment}
\end{equation}

For a mass-weighted objective:
\begin{equation}
  M^{\mathrm{mass}}
  =\sum_{q\in\Place}
  w_{i(q)}
  \left(\frac{g(z_q+d_q^z/2)}{1000}\right)x_q.
  \label{eq:mass-moment}
\end{equation}

The factor $z_q+d_q^z/2$ is the vertical center of the selected box in grid
units.  Multiplication by $g$ converts it to millimetres, and division by
$1000$ converts millimetres to metres.  The density or mass coefficient is
fixed data, so both expressions are linear in $x$.

\section{Multi-objective solution strategy}

The implementation solves a sequence of MILPs rather than combining all
objectives into one arbitrary weighted sum.  For each feasible pallet count
$K$, it performs the following lexicographic/Pareto procedure:

\begin{enumerate}[label=\arabic*.]
  \item Minimize the height spread $s$.
  \item Add a constraint fixing $s$ to its optimum and minimize the
    accessibility count $A$.
  \item Retain the best accessibility value as a starting point, then minimize
    the selected vertical moment $M$.
  \item For successive accessibility bounds $\varepsilon$, add
    $A\le\varepsilon$ and minimize $M$ again.  The resulting nondominated
    solutions form the reported discrete Pareto front.
\end{enumerate}

Every optimization problem in this sequence has the same binary placement
variables and core feasibility constraints.  Only the objective and the
additional Pareto bound change between stages.

\section{Summary of symbol origins}

\begin{table}[h]
\centering
\small
\begin{tabularx}{\linewidth}{@{}lXX@{}}
\toprule
Symbol & Meaning & Origin \\
\midrule
$\Item,\ d_i,\ w_i,\rho_i,\Prio_i$ & Item data & Input instance JSON \\
$L,W,H$ & Pallet dimensions & Input instance JSON, divided by grid size \\
$\Pallet$ & Candidate pallet indices & Solver configuration \\
$\Place$ & Feasible placements & Orientation/coordinate preprocessing \\
$\Cell$ & Discrete occupied cells & Placement geometry \\
$A_q,A_{qr}$ & Base and overlap areas & Placement geometry \\
$f_i,c_i,\varphi_i$ & Handling attributes & Input instance metadata \\
$K,\alpha$ & Pallet count and support fraction & Solve loop/configuration \\
$x_q,y_p,h_p,H^{\max},H^{\min},s,a_{ij}$ & Decision variables & Created by the MILP \\
\bottomrule
\end{tabularx}
\caption{Origin of the notation used in the abstract formulation.}
\label{tab:symbol-origins}
\end{table}

\clearpage
\section{Complete compact MILP}

For reference, the entire abstract model is collected below.  Let
$\mu_i=\rho_i$ for the density-weighted objective and $\mu_i=w_i$ for the
mass-weighted objective.  Define
\begin{align}
  A(x,a) &:= \sum_{(i,j)\in\mathcal{A}} a_{ij},
  \label{eq:complete-A}\\
  M(x) &:= \sum_{q\in\Place}
  \mu_{i(q)}\frac{g(z_q+d_q^z/2)}{1000}x_q.
  \label{eq:complete-M}
\end{align}
For each fixed pallet count $K$, the abstract multi-objective MILP is

\begin{subequations}
\label{model:complete}
\begin{alignat}{3}
  \underset{x,y,h,H^{\max},H^{\min},s,a}{\operatorname{Pareto-minimize}}
  \quad & \bigl(s,\ A(x,a),\ M(x)\bigr)
  \label{model:objective}\\[0.3em]
  \text{subject to}\quad
  & \sum_{q\in\Place_i}x_q = 1
  && \qquad \forall i\in\Item,
  \label{model:place-once}\\
  & \sum_{q\in\Place_p}x_q \ge y_p
  && \qquad \forall p\in\Pallet,
  \label{model:pallet-nonempty}\\
  & \sum_{q\in\Place_{ip}}x_q \le y_p
  && \qquad \forall i\in\Item,\ p\in\Pallet,
  \label{model:item-pallet-link}\\
  & \sum_{p\in\Pallet}y_p = K,
  &&
  \label{model:exact-pallets}\\
  & y_p \ge y_{p+1}
  && \qquad p=1,\ldots,P^{\max}-1,
  \label{model:pallet-symmetry}\\
  & h_{p(q)} \ge t_qx_q
  && \qquad \forall q\in\Place,
  \label{model:placement-height}\\
  & h_p \le Hy_p
  && \qquad \forall p\in\Pallet,
  \label{model:unused-height}\\
  & H^{\max}\ge h_p
  && \qquad \forall p\in\Pallet,
  \label{model:max-height}\\
  & H^{\min}\le h_p+H(1-y_p)
  && \qquad \forall p\in\Pallet,
  \label{model:min-height}\\
  & s=H^{\max}-H^{\min},
  &&
  \label{model:spread}\\
  & \sum_{q\in\Place(c)}x_q\le1
  && \qquad \forall c\in\Cell,
  \label{model:cell-clique}\\
  & \sum_{r\in\Support(q)}A_{qr}x_r
      \ge \alpha A_qx_q
  && \qquad \forall q\in\Place:\ z_q>0,
  \label{model:support}\\
  & \sum_{q\in B}x_q+\sum_{q\in U}x_q\le1
  && \qquad \forall(B,U)\in\mathcal{H},
  \label{model:hard-relations}\\
  & a_{ij}\ge
      \sum_{q\in B}x_q+\sum_{q\in U}x_q-1
  && \qquad \substack{\forall(i,j)\in\mathcal{A},\\
                       \forall(B,U)\in\mathcal{R}_{ij}},
  \label{model:access-relations}\\
  & \sum_{q\in\Place_i}r_{iq}x_q
      \le \sum_{q\in\Place_j}r_{jq}x_q
  && \qquad \substack{\forall\text{ consecutive identical}\\
                       \text{items }i,j\text{ in a group }G\in\mathcal{G}},
  \label{model:identical-symmetry}\\[0.2em]
  & x_q\in\{0,1\}
  && \qquad \forall q\in\Place,
  \label{model:x-domain}\\
  & y_p\in\{0,1\},\quad 0\le h_p\le H
  && \qquad \forall p\in\Pallet,
  \label{model:yh-domain}\\
  & a_{ij}\in\{0,1\}
  && \qquad \forall(i,j)\in\mathcal{A},
  \label{model:a-domain}\\
  & 0\le H^{\min},H^{\max},s\le H.
  &&
  \label{model:height-domains}
\end{alignat}
\end{subequations}

The implementation realizes the Pareto objective in
\eqref{model:objective} through repeated scalar MILPs.  It first minimizes
$s$, fixes the optimal spread, and minimizes $A$.  It then solves
$\min M(x)$ under bounds $A(x,a)\le\varepsilon$ for a sequence of integer
values $\varepsilon$.  The outer loop repeats this process for admissible
values of $K$.  Consequently, the compact model above and the implemented
sequence describe the same feasible packing decisions and objective trade-offs;
the latter is the computational procedure used to enumerate nondominated
solutions.

\end{document}
